\chapter{Relations and Orders}
\label{ch:relations-orders}

This chapter covers equivalence relations, partial orders,
transitivity, monotonicity, and a light introduction to
well-foundedness. It is also the first chapter that recasts
definitions already used informally elsewhere — extension and
constraint on a history of witnessed facts — as instances of
ordinary relational and order-theoretic structure, rather than
leaving them as isolated metaphors.

\section{Equivalence Relations}

An equivalence relation bundles three properties: reflexivity
($x \sim x$ for every $x$), symmetry ($x \sim y \to y \sim x$), and
transitivity ($x \sim y \to y \sim z \to x \sim z$). All three
matter independently — the countermodel below exhibits a relation
with the first two properties and not the third.

\begin{experiment}{Same parity is an equivalence relation}
\begin{lean}
def SameParity (x y : Fin 4) : Bool := x.val % 2 == y.val % 2

example : ∀ x : Fin 4, SameParity x x = true := by decide
example : ∀ x y : Fin 4, SameParity x y = SameParity y x := by decide
example : ∀ x y z : Fin 4,
    SameParity x y = true → SameParity y z = true →
    SameParity x z = true := by decide
\end{lean}
All three checks are finite and decidable over \texttt{Fin 4}, so
each closes by exhaustive search rather than by a hand-written
argument — reflexivity and symmetry are visibly built into the
definition of \texttt{SameParity} itself, and transitivity, though
less obvious by inspection, is confirmed the same mechanical way.
\end{experiment}

\begin{countermodel}{Reflexive and symmetric does not imply transitive}
\begin{lean}
def Close (x y : Fin 3) : Bool :=
  x == y || (x == 0 && y == 1) || (x == 1 && y == 0) ||
  (x == 1 && y == 2) || (x == 2 && y == 1)

example : ∀ x : Fin 3, Close x x = true := by decide
example : ∀ x y : Fin 3, Close x y = Close y x := by decide

example : ¬ (∀ x y z : Fin 3,
    Close x y = true → Close y z = true → Close x z = true) := by
  decide
\end{lean}
\texttt{Close} relates $0$ to $1$ and $1$ to $2$, but not $0$ to
$2$ — reflexive and symmetric by construction, and the failure of
transitivity is exactly the familiar fact that ``close enough''
relations need not chain: two short hops can add up to one that
is not short.
\end{countermodel}

\section{Partial Orders and Monotonicity}

\begin{experiment}{Extension as a partial order}
A history's set of witnessed facts only grows: extension is
reflexive and transitive.
\begin{lean}
structure History (alpha : Type) where
  witnessed : Set alpha

def Extends {alpha : Type} (earlier later : History alpha) : Prop :=
  earlier.witnessed ⊆ later.witnessed

theorem extends_refl {alpha : Type} (h : History alpha) :
    Extends h h := fun x hx => hx

theorem extends_trans {alpha : Type} {h1 h2 h3 : History alpha}
    (h12 : Extends h1 h2) (h23 : Extends h2 h3) :
    Extends h1 h3 := fun x hx => h23 (h12 hx)
\end{lean}
Reflexivity and transitivity together make \texttt{Extends} a
preorder on histories.
\end{experiment}

For \texttt{History} exactly as defined here
— a bare wrapper around one \texttt{Set} field — antisymmetry
actually does hold: two histories extending each other have equal
witnessed sets by set extensionality, and since \texttt{witnessed}
is the only field, equal witnessed sets means equal histories. The
weaker claim of ``preorder'' is stated deliberately anyway, because
it is the property that survives a natural enrichment this bare
structure invites — a \texttt{History} carrying additional fields
(a timestamp, a provenance tag, a source identifier) could have two
genuinely distinct instances with identical witnessed sets, which
would extend each other in both directions without being equal.
Stating only what is guaranteed to keep holding is worth the
weaker-sounding claim.

\begin{experiment}{Constraint narrows admissibility, never expands it}
\begin{lean}
def Constrain {alpha : Type} (admissible constraint : Set alpha) :
    Set alpha := admissible ∩ constraint

theorem constrain_subset {alpha : Type}
    (admissible constraint : Set alpha) :
    Constrain admissible constraint ⊆ admissible := fun x hx => hx.1
\end{lean}
This is almost too simple to look like a theorem, and that is
exactly the point made concrete: ``a constraint can only narrow what
is admissible, never expand it'' is not asserted anywhere in this
framework as a slogan — it is this one line, checked the same way
every other claim in this book has been.
\end{experiment}

\section{Well-Foundedness}

A relation $\prec$ is \emph{well-founded} if there is no infinite
descending chain $\cdots \prec x_2 \prec x_1 \prec x_0$ — every
descent, however it is chosen, must eventually stop. This is the
fact underneath two things already taken for granted in this book:
structural recursion, and the assumption that a constraint process
cannot run forever.

Lean accepts a recursive definition without an explicit termination
proof exactly when each recursive call is made on an argument
strictly smaller under a well-founded relation, ordinarily
\texttt{Nat}'s $<$ or a type's own notion of being a smaller
subterm.
\begin{lean}
def countdown : Nat → Nat
  | 0 => 0
  | n + 1 => countdown n
\end{lean}
The recursive call \texttt{countdown n} is made on \texttt{n}, which
is structurally smaller than \texttt{n + 1} — the termination
checker verifies precisely this, silently, and the well-foundedness
of $<$ on \texttt{Nat} is exactly what rules out an infinite chain
of ever-smaller natural numbers that could make the check fail.

The same fact underwrites constraint termination directly: if a
process constrains its admissible set at every step and some
$\texttt{Nat}$-valued measure (the size of a finite admissible set,
say) strictly decreases at each constraining step, the process
cannot continue indefinitely — an infinite run would be exactly the
infinite descending chain well-foundedness rules out. Nothing about
constraint or admissibility needs to be assumed further; termination
follows from the well-foundedness already built into $<$ on the
natural numbers.

\section{Exercises}

\begin{exercise}
Prove that intersecting an admissible set with a constraint set never
enlarges it: $\mathrm{Constrain}(A, C) \subseteq A$.
\end{exercise}

\begin{exercise}
Prove that constraining twice by the same set changes nothing further:
$\mathrm{Constrain}(\mathrm{Constrain}(A, C), C) = \mathrm{Constrain}(A, C)$.
\end{exercise}
